\chapter*{Notation and Conventions}
\addcontentsline{toc}{chapter}{Notation and Conventions}

This manuscript distinguishes the \emph{true plant state} from the
\emph{telemetry-delivered state} at every stage of the argument.  The symbol
$x_t$ denotes the latent physical state, while $\hat{x}_t$ and $z_t$ denote the
controller-visible observation surface after parsing, delay, or fault
processing.  The nominal intelligence proposes a candidate action
$a_t^{\mathrm{cand}}$; ORIUS releases either the repaired action
$a_t^{\mathrm{safe}}$ or an admissible fallback action.

The reliability score $w_t \in [0,1]$ is a runtime trust signal, not a
probability by definition.  Whenever a theorem uses $w_t$ inside a probabilistic
bound, that move is justified by an explicit theorem-local contract.  The
observation-consistent state set is written as $U_t$ in the universal chapters
and as the domain-instantiated conformal tube in the reference proofs.  The
safe set is written as $\mathcal{S}$ or, in the battery witness, through the
SOC interval constraint.

Formal counts in this dissertation are computed from the active master-manuscript
include tree, not from every standalone \texttt{.tex} file in the repository.
This avoids double-counting archived chapter surfaces and keeps the published
formal inventory aligned to the canonical PDF.  Appendix~\ref{app:notation} and
Chapter~\ref{ch:assumptions} remain the authoritative sources for the complete
notation table and the assumption register.
